\documentclass[../../main.tex]{subfiles}

\begin{document}

\chapter{Otentikasi dan Tanda Tangan Digital}

\begin{learningoutcome}
Setelah mempelajari bab ini, mahasiswa diharapkan mampu:
\begin{itemize}
    \item Menjelaskan konsep otentikasi dan pentingnya non-repudiation dalam komunikasi data (Sub-CPMK 1.4).
    \item Menguraikan mekanisme kerja tanda tangan digital menggunakan kriptografi kunci-publik (Sub-CPMK 1.4).
    \item Menganalisis perbedaan antara tanda tangan digital berbasis RSA dan DSS (Digital Signature Standard) (Sub-CPMK 1.4).
    \item Menjelaskan peran Public Key Infrastructure (PKI) dan sertifikat digital X.509 dalam mendistribusikan kunci publik secara aman (Sub-CPMK 1.4).
    \item Mengevaluasi alur kerja protokol SSL/TLS dalam mengamankan sesi komunikasi internet (Sub-CPMK 1.4).
\end{itemize}
\end{learningoutcome}

\subfile{section-01}
\subfile{section-02}
\subfile{section-03}
\subfile{section-04}
\section{Contoh Terhitung}
Setiap langkah berikut diverifikasi secara numerik dengan Python.

\subsection{Tanda Tangan RSA dengan Bilangan Kecil}
Pilih $p = 61$, $q = 53$, sehingga
\[ n = p\,q = 3233, \qquad \varphi(n) = (p-1)(q-1) = 60 \cdot 52 = 3120 . \]
Ambil $e = 17$ (relatif prima terhadap $3120$), maka $d = e^{-1} \bmod \varphi(n) = 2753$ karena $17 \cdot 2753 = 46801 = 15 \cdot 3120 + 1$. Pasangan kunci: privat $d = 2753$, publik $(e,n) = (17, 3233)$.

Misalkan fungsi hash mainan untuk pesan 1 karakter \code{'A'} memberi $h = H(m) = 65$ (kode ASCII-nya, dan $h < n$). Penandatanganan:
\[ s = h^{d} \bmod n = 65^{2753} \bmod 3233 = 588 . \]
Verifikasi oleh penerima:
\[ s^{e} \bmod n = 588^{17} \bmod 3233 = 65 = h \;\checkmark \]
Tanda tangan valid. Sekarang uji pemalsuan: bila penyerang mengubah pesan sehingga hash mainannya menjadi $h' = 66$, penerima menghitung $h' = 66$ sedangkan dari tanda tangan lama diperoleh $s^{e} \bmod n = 65 \neq 66$, sehingga tanda tangan ditolak.

\subsection{Tanda Tangan DSA dengan Parameter Kecil}
Ambil prima $p = 23$ dan $q = 11$ (karena $11 \mid 22 = p-1$), generator $g = 2$; cek ordonya: $2^{11} \bmod 23 = 2048 \bmod 23 = 1$, jadi $g$ berorde $q$. Kunci privat $x = 5$ memberi kunci publik
\[ y = g^{x} \bmod p = 2^{5} \bmod 23 = 32 \bmod 23 = 9 . \]
Untuk hash mainan $h = 10$ dan bilangan acak sesi $k = 3$ (dengan $k^{-1} \bmod 11 = 4$):
\[ r = (g^{k} \bmod p) \bmod q = (2^{3} \bmod 23) \bmod 11 = 8 \bmod 11 = 8, \]
\[ s = k^{-1}(h + x\,r) \bmod q = 4\,(10 + 5 \cdot 8) \bmod 11 = 4 \cdot 50 \bmod 11 = 200 \bmod 11 = 2 . \]
Pasangan tanda tangan $(r, s) = (8, 2)$.

Verifikasi: $w = s^{-1} \bmod 11 = 2^{-1} \bmod 11 = 6$,
\[ u_1 = h\,w \bmod 11 = 10 \cdot 6 \bmod 11 = 5, \qquad u_2 = r\,w \bmod 11 = 8 \cdot 6 \bmod 11 = 4, \]
\[ v = \bigl(g^{u_1} y^{u_2} \bmod p\bigr) \bmod q = \bigl(2^{5} \cdot 9^{4} \bmod 23\bigr) \bmod 11 = (9 \cdot 6 \bmod 23) \bmod 11 = 8 . \]
Karena $v = 8 = r$, tanda tangan \textbf{valid}.

\subsection{Dampak Pemakaian Ulang \texorpdfstring{$k$}{k}}
Jika dua pesan dengan hash berbeda $h_1, h_2$ ditandatangani memakai $k$ yang sama, maka $r$ keduanya sama dan penyerang dapat menyelesaikan
\[ x = \frac{s_1 - s_2}{r} \cdot \frac{1}{h_1 - h_2} \pmod q, \]
sehingga seluruh kunci privat bocor. Inilah sebabnya ECDSA/DSA menuntut $k$ acak yang unik per tanda tangan \citep{stallings2017}.

\begin{obeactivity}
\textbf{Aktivitas 16.1: Analisis Sertifikat Digital Browser}
1. Buka browser (Chrome/Firefox) dan akses situs dengan HTTPS (misal: google.com).
2. Klik ikon gembok di samping URL $\rightarrow$ \textit{Connection is secure} $\rightarrow$ \textit{Certificate is valid}.
3. Identifikasi siapa penerbit sertifikatnya (CA), kapan masa berlakunya, dan apa kunci publik yang digunakan.
4. Diskusikan apa yang terjadi jika sertifikat tersebut kedaluwarsa atau tidak diterbitkan oleh CA terpercaya.
\end{obeactivity}



\begin{obereflection}
\textbf{Latihan 16.1}
1. Mengapa tanda tangan digital menggunakan hash dari pesan, bukan mengenkripsi seluruh pesan dengan kunci privat?
2. Apa perbedaan antara otentikasi menggunakan password dibandingkan dengan otentikasi menggunakan tanda tangan digital?
3. Bagaimana PKI mencegah serangan \textit{Man-in-the-Middle} dalam distribusi kunci publik?

\textbf{Latihan 16.2}
4. Jelaskan mengapa MAC tidak dapat dipakai untuk mencapai nirpenyangkalan, padahal MAC juga membuktikan keaslian pesan. Kaitkan jawaban Anda dengan pertanyaan ``siapa yang memegang kunci''.
5. Sebuah protokol membuktikan identitas pengguna dengan mengirimkan nilai hash dari kata sandinya. Penyerang yang menyadap pertukaran itu tidak mengetahui kata sandinya, tetapi ia tetap dapat menyamar. Serangan apa yang ia lakukan, dan bagaimana \textit{nonce} menutupnya?
6. Pada contoh DSA bab ini, penandatangan memakai $k = 3$ untuk dua pesan. Hitunglah sendiri $s_2$ untuk $h_2 = 7$, lalu ulangi pemulihan $k$ dan $x$. Periksa bahwa hasilnya sama dengan yang ada di bab.

\textbf{Latihan 16.3}
7. Dengan parameter mainan $n = 3233$ dan $d = 2753$ pada bab ini, tanda tangan untuk $h_1 = 65$ adalah $588$ dan untuk $h_2 = 66$ adalah $2215$. Tunjukkan bahwa $588 \cdot 2215 \bmod 3233$ adalah tanda tangan yang sah atas hash $1057$, dan jelaskan mengapa serangan ini tidak akan berhasil bila skema padding PSS dipakai.
8. Kunci publik Ed25519 hanya 32 byte, sedangkan kunci publik ECDSA P-256 adalah 65 byte, padahal keduanya sama-sama berada di atas kurva 256 bit. Jelaskan sumber perbedaan itu.
9. Mengapa penyerang tidak dapat menyamar hanya dengan menyalin sertifikat sebuah situs web? Sebutkan pesan dalam jabat tangan TLS yang menutup celah tersebut.

\textbf{Latihan 16.4}
10. Pada Tabel~\ref{tab:jaminan-tls}, tahap \textit{ServerHello} belum memberikan jaminan keaslian pengirim. Mengapa? Dan apa yang membuat tahap berikutnya berbeda?
11. Rancang sebuah skenario di mana rantai sertifikat yang seluruhnya sah tetap menghasilkan koneksi yang tidak aman. Petunjuk: pikirkan tentang ekstensi \textit{Key Usage} dan tentang panjang kunci yang sudah usang.
12. Sebuah organisasi hendak menerbitkan sertifikat untuk perangkat \textit{Internet of Things} berdaya rendah yang harus menandatangani pesan setiap detik. Algoritma mana yang Anda pilih di antara empat pada Tabel~\ref{tab:perbandingan-tanda-tangan}, dan mengapa? Sebutkan juga kerugian dari pilihan Anda.
\end{obereflection}

\section{Rangkuman}
\begin{itemize}
    \item Otentikasi entitas membuktikan identitas pihak pada saat sesi dibangun (dengan \textit{nonce} agar tahan \textit{replay}), sedangkan otentikasi pesan membuktikan keaslian dan keutuhan tiap pesan.
    \item Tanda tangan digital mengikat identitas pengirim pada isi pesan melalui hash yang ditandatangani dengan kunci privat; verifikasinya memakai kunci publik.
    \item RSA menandatangani dengan $s = H(m)^d \bmod n$ dan diverifikasi dengan $H(m) = s^e \bmod n$; DSS memakai DSA/ECDSA pada subgrup prima dengan tanda tangan $(r,s)$.
    \item Pesan selalu di-hash lebih dulu karena hashing memampatkan pesan, mengefisienkan komputasi, dan mengikat tanda tangan pada isi pesan; keamanannya bergantung pada ketahanan kolisi ($2^{n/2}$).
    \item \textit{Non-repudiation}, yaitu jaminan pengirim tidak dapat menyangkal pesannya, hanya dapat dijamin oleh tanda tangan digital, bukan oleh MAC yang bersifat simetris.
    \item Sertifikat X.509 mengikat identitas subjek dengan kunci publiknya dan ditandatangani CA; keabsahannya diperiksa melalui rantai kepercayaan Root CA $\to$ Intermediate CA $\to$ sertifikat server.
    \item Jabat tangan TLS memakai kriptografi asimetris untuk otentikasi dan penyepakatan kunci, lalu kriptografi simetris untuk sesi; TLS 1.3 mempersingkat proses dan mewajibkan \textit{forward secrecy}.
\end{itemize}

\begin{obeassessment}
\textbf{Kuis Bab 16}
1. Jelaskan kaitan antara hash, kunci privat, dan kunci publik dalam proses tanda tangan digital.
2. Apa peran sertifikat X.509 dalam ekosistem Public Key Infrastructure?
3. Mengapa protokol TLS menggunakan kombinasi kriptografi asimetris (di awal) dan simetris (selama sesi)?
\end{obeassessment}
\begin{competencychecklist}
\begin{tabularx}{\linewidth}{|X|c|}
\hline
Indikator Kompetensi & Tercapai \\
\hline
Saya dapat membedakan otentikasi dan non-repudiation & [ ] \\
\hline
Saya dapat menguraikan alur pembuatan dan verifikasi tanda tangan digital & [ ] \\
\hline
Saya dapat menjelaskan fungsi CA dan sertifikat X.509 & [ ] \\
\hline
Saya dapat menganalisis mekanisme handshake SSL/TLS & [ ] \\
\hline
Saya dapat mengidentifikasi komponen PKI dalam aplikasi HTTPS & [ ] \\
\hline
\end{tabularx}
\end{competencychecklist}


\end{document}
